\subsection*{NT-H12. Class-group torsion from the field's own Dedekind RH}
Fix integers $n\ge2$ and $\ell\ge2$. For every degree-$n$ number field
$K$ with absolute discriminant $D_K$ whose Dedekind zeta function satisfies
RH, and every $\epsilon>0$,
\[
 |\operatorname{Cl}_K[\ell]|
       \ll_{n,\ell,\epsilon}
 D_K^{\,1/2-1/(2\ell(n-1))+\epsilon}.                   \tag{NT-H12a}
\]
Only RH for $\zeta_K$ is needed. In particular $M$ supplies the hypothesis
conditionally, and adopted $A$ supplies it when $K/\mathbb Q$ is abelian.
No RH for a normal closure or Hilbert class field is required by the
argument below. The exponent is the classical Ellenberg--Venkatesh
bound, not a new class-group exponent. Our point is the exact
Dedekind-only input reduction.

The precise cited algebraic input is Ellenberg--Venkatesh,
\emph{Reflection principles and bounds for class group torsion},
author preprint, Lemma3 and its proof, p.6: if
$0<\delta<1/(2\ell(n-1))$ and there are $N$ prime ideals of $K$
of norm at most $D_K^\delta$, unramified over $\mathbb Q$ and not
extensions of prime ideals of any proper intermediate field, then
\[
 |\operatorname{Cl}_K[\ell]|
       \ll_{n,\ell,\epsilon,\delta}D_K^{1/2+\epsilon}/N.
                                                        \tag{NT-H12b}
\]
The source's Proposition1, p.7, applies this lemma using completely
split primes. We instead count degree-one prime ideals. The implicit
constant in (NT-H12a) is not assigned an unchecked numerical value.

\begin{proof}
Let $\pi_K(x)$ count all prime ideals of $K$ of norm at most $x$.
Apply NT-C02 to the trivial Galois extension $K/K$. Its conjugacy density
is one, there are no relative ramified primes to omit, and RH for
$\zeta_K$ gives, for $x\ge2$,
\[
 |\pi_K(x)-\operatorname{Li}_2(x)|\le B_K(x),
\]
where the completely explicit error is
\[
 B_K(x)=\sqrt x\left[
 \left(\frac1{2\pi}+\frac3{\log x}\right)\log D_K+
 \left(\frac{\log x}{8\pi}+\frac1{4\pi}+\frac6{\log x}\right)n
                  \right].                            \tag{NT-H12c}
\]
Every prime ideal of residue degree at least two above its rational
prime has underlying prime at most $\sqrt x$. There are at most $n$
prime ideals above each rational prime, so such ideals number at most
$n\pi(\sqrt x)\le n\sqrt x$. There are at most $n\omega(D_K)$
prime ideals above rational primes ramified in $K$. Therefore the
number $N_K^{(1),\mathrm{ur}}(x)$ of unramified degree-one prime ideals
satisfies the useful finite inequality
\[
 N_K^{(1),\mathrm{ur}}(x)
   \ge\operatorname{Li}_2(x)-B_K(x)-n\sqrt x-n\omega(D_K).
                                                        \tag{NT-H12d}
\]
Here $\omega$ counts distinct rational prime divisors.

These degree-one ideals meet the remaining condition in (NT-H12b).
Indeed, if a prime ideal $\mathfrak p$ of norm $p$ were
$\mathfrak q\mathcal O_K$ for a prime ideal of a proper subfield $F$,
then multiplicativity of the norm under ideal extension would give
\[
 p=N_K(\mathfrak p)=(N_F\mathfrak q)^{[K:F]}.
\]
The right side is a proper positive integral power of an integer
at least two, and cannot be the rational prime $p$.

Fix $\delta>0$ and put $x=D_K^\delta$. For large $D_K$, the bound
$\operatorname{Li}_2(x)\ge x/(2\log x)$ follows by integrating on
$[x/2,x]$. The error in (NT-H12c) is
$O_{n,\delta}(D_K^{\delta/2}\log D_K)$, and
$\omega(D_K)\le\log D_K/\log2$. Consequently
\[
 N_K^{(1),\mathrm{ur}}(D_K^\delta)
                  \gg_{n,\delta}D_K^\delta/\log D_K.   \tag{NT-H12e}
\]
Write $\alpha=1/(2\ell(n-1))$. It suffices to prove (NT-H12a) for
$0<\epsilon<\alpha$, since a result with any smaller positive epsilon
implies the assertion with a larger one. Choose
$\delta=\alpha-\epsilon/3$ and apply (NT-H12b) with its epsilon
parameter equal to $\epsilon/3$. Equations (NT-H12b) and (NT-H12e)
give
\[
 |\operatorname{Cl}_K[\ell]|
   \ll D_K^{1/2-\alpha+2\epsilon/3}\log D_K
   \ll_{n,\ell,\epsilon}D_K^{1/2-\alpha+\epsilon}.
\]
The bounded-discriminant cases are absorbed into the constant by
Hermite's finiteness theorem (Milne, \emph{Algebraic Number Theory}
v3.08, Theorem8.43, printed pp.151--152). This proves the claim.
\end{proof}
The geometry is the torsion subgroup of
$\operatorname{Pic}(\operatorname{Spec}\mathcal O_K)$.
The source's Arakelov argument incorporates all archimedean places;
our analytic step counts rational degree-one closed points without
requiring a completely split Frobenius class.

\subsection*{NT-H13. Bounds for unramified cyclic covers and Kummer torsors}
Retain $n\ge2$, $K$, and RH for $\zeta_K$ from NT-H12.
For a prime $\ell$, the number of $K$-isomorphism classes of cyclic
extensions $E/K$ of degree $\ell$, unramified at every finite place
and split at every real place, is exactly
\[
 c_\ell(K)=\frac{|\operatorname{Cl}_K[\ell]|-1}{\ell-1}
       \ll_{n,\ell,\epsilon}
       D_K^{\,1/2-1/(2\ell(n-1))+\epsilon}.             \tag{NT-H13a}
\]
For odd $\ell$ the real-place condition is automatic. For $\ell=2$
it is essential to this ordinary-class-group formula.

For any integer $\ell\ge2$, let $S$ be the set of primes of $K$
above divisors of $\ell$, let $U=\operatorname{Spec}\mathcal O_K[1/\ell]$,
and let $w_K$ be the number of roots of unity in $K$. The number of
isomorphism classes of $\mu_\ell$-torsors over $U$ is
\[
 \begin{split}
 |H^1_{\text{ét}}(U,\mu_\ell)|
  &=\ell^{r_1+r_2-1+|S|}\gcd(\ell,w_K)
                             |\operatorname{Pic}(U)[\ell]|\\
  &\le\ell^{r_1+r_2+|S|}|\operatorname{Cl}_K[\ell]|\\
  &\ll_{n,\ell,\epsilon}
                 D_K^{\,1/2-1/(2\ell(n-1))+\epsilon}.
 \end{split}                                           \tag{NT-H13b}
\]
The finite coefficient $\ell^{r_1+r_2+|S|}$ is at most
$\ell^{n(1+\omega(\ell))}$. These torsors need not be connected;
$\mu_\ell$ need not be a constant group scheme over $K$.

\begin{proof}
Global class field theory identifies the ordinary Hilbert class field
Galois group with $\operatorname{Cl}_K$, with real places split.
The convention is stated explicitly in Milne, \emph{Algebraic Number
Theory} v3.08, Remark4.11 and its footnote, printed p.72; the general
reciprocity/existence input is Milne, \emph{Class Field Theory}, V3.5--3.6.
Degree-$\ell$ cyclic subextensions correspond to kernels of surjective
homomorphisms $\operatorname{Cl}_K\to\mathbb Z/\ell\mathbb Z$.
For prime $\ell$ every nonzero homomorphism is surjective, and
\[
 |\operatorname{Hom}(\operatorname{Cl}_K,\mathbb Z/\ell\mathbb Z)|
      =|\operatorname{Cl}_K[\ell]|.
\]
Exactly $\ell-1$ such homomorphisms have each kernel, by postcomposition
with automorphisms of the target. This proves the equality in
(NT-H13a), and NT-H12 gives the bound. A real-place inertia group has
order at most two, so it is trivial in an extension of odd prime degree.

Put $E_{K,S}=\mathcal O_K[1/\ell]^\times$. Since $\ell$ is invertible
on $U$, the étale Kummer sequence and
$H^1_{\text{ét}}(U,\mathbb G_m)=\operatorname{Pic}(U)$ give
\[
 0\longrightarrow
 \mathcal O_K[1/\ell]^\times/igl(\mathcal O_K[1/\ell]^\times\bigr)^\ell
 \longrightarrow H^1_{\text{ét}}(U,\mu_\ell)
 \longrightarrow\operatorname{Pic}(U)[\ell]
 \longrightarrow0.                                    \tag{NT-H13c}
\]
A precise source is Stacks Project, Section59.28, Tag03PK,
Lemma59.28.1 and its following cohomology sequence; Lemma59.28.5,
Tag040Q, identifies the classes as line bundles with an $\ell$th-power
trivialization. Thus this cohomology group is not being confused with
Picard torsion alone.

The $S$-unit theorem gives
$\mathcal O_K[1/\ell]^\times\simeq\mu(K)\times
\mathbb Z^{r_1+r_2-1+|S|}$ as abelian groups. Here $\mu(K)$ is cyclic
of order $w_K$. See Milne, \emph{Algebraic Number Theory} v3.08,
Theorem5.11 and its complete proof, printed p.90. Taking the quotient
by $\ell$th powers has cardinality
$\ell^{r_1+r_2-1+|S|}\gcd(\ell,w_K)$, proving the first equality in
(NT-H13b).

Localization makes $\operatorname{Pic}(U)$ the quotient of
$\operatorname{Cl}_K$ by the classes of primes in $S$. For any
surjection $A\to B$ of finite abelian groups, the map
$A/\ell A\to B/\ell B$ is surjective and
$|A[\ell]|=|A/\ell A|$, $|B[\ell]|=|B/\ell B|$.
It follows that $|\operatorname{Pic}(U)[\ell]|
\le|\operatorname{Cl}_K[\ell]|$; no injection on torsion induced by
localization is asserted. Apply NT-H12 and
$\gcd(\ell,w_K)\le\ell$, $r_1+r_2\le n$,
$|S|\le n\omega(\ell)$ to finish the proof.
\end{proof}
This is a literal finite étale/Kummer interpretation of a conditional
class-group bound: cyclic quotients of the arithmetic fundamental
group with the stated real conditions, and torsors under the finite
étale group scheme $\mu_\ell$ over $U$. It gives no bound for arbitrary
nonabelian covers or for all Selmer groups without further reduction.
